% sec13_8.tex — Section 13.8: Solving the Wave Equation Using Laplace Transforms (with Complex Variables)
\section{Solving the Wave Equation Using Laplace Transforms (with Complex Variables)}
\label{sec:13.8}

In Sec.~13.6 we showed that
\begin{align}
\text{PDE:}\quad &
\begin{array}{|c|}
\hline
\displaystyle\pdd{u}{t} = c^2\pdd{u}{x} \\[4pt]
\hline
\end{array}
\label{eq:13.8.1} \\[8pt]
\text{BC:}\quad &
\begin{array}{|c|}
\hline
u(0,t) = 0 \\
u(L,t) = b(t) \\
\hline
\end{array}
\label{eq:13.8.2} \\[8pt]
\text{IC:}\quad &
\begin{array}{|c|}
\hline
u(x,0) = 0 \\[2pt]
\displaystyle\pd{u}{t}(x,0) = 0 \\
\hline
\end{array}
\label{eq:13.8.3}
\end{align}
could be analyzed by introducing $\overline{U}(x,s)$, the Laplace transform in $t$ of $u(x,t)$.  We obtained $u(x,t)$,
\begin{equation}
u(x,t) = -c^2\int_0^t b(t_0)\,\pd{G}{x_0}(x,t;L,t_0)\,dt_0,
\label{eq:13.8.4}
\end{equation}
in terms of the Green's function.  The Laplace transform of this influence function was known [see \eqref{eq:13.6.15}]:
\begin{equation}
\pd{\overline{G}}{x_0}(x,s;L,0) = F(s) = -\frac{\sinh(x/c)s}{c^2\sinh(L/c)s}.
\label{eq:13.8.5}
\end{equation}

In this section we use the complex inversion integral for the Laplace transform:
\begin{equation}
\pd{G}{x_0}(x,t;L,0) = \frac{1}{2\pi i}\int_{\gamma-i\infty}^{\gamma+i\infty}
-\frac{\sinh(x/c)s}{c^2\sinh(L/c)s}\,e^{st}\ds.
\label{eq:13.8.6}
\end{equation}
The singularities of $F(s)$ only are simple poles $s_n$, located at the zeros of the denominator:
\begin{equation}
\sinh\frac{L}{c}\,s_n = 0.
\label{eq:13.8.7}
\end{equation}
However, $s = 0$ is not a pole since, near $s = 0$, $F(s) \approx -(x/c)s/[c^2(L/c)s] \ne \infty$.  There is an infinite number of these poles located on the imaginary axis:
\begin{equation}
\frac{L}{c}\,s_n = in\pi, \qquad n = \pm 1, \pm 2, \pm 3, \dots.
\label{eq:13.8.8}
\end{equation}
The location of the poles, $s_n = ic(n\pi/L)$, corresponds to the eigenvalues, $\lambda = (n\pi/L)^2$.  This follows from the singularity property of Laplace transforms (see Sec.~13.2).  The singularity of the Laplace transform [$s = ic(n\pi/L)$] corresponds to a complex exponential solution [$e^{ic(n\pi/L)t}$].

The residue at each pole may be evaluated:
\[
\operatorname{res}(s_n) = \frac{R(s_n)}{Q'(s_n)}
= \frac{-\sinh(x/c)s_n\,e^{s_n t}}{cL\cosh(L/c)s_n}
= \frac{-i\sin n\pi x/L\,e^{i(n\pi ct/L)}}{cL\cos n\pi},
\]
since $\sinh ix = i\sin x$ and $\cosh ix = \cos x$.  Thus, the influence function for this problem is
\begin{align}
\frac{\partial}{\partial x_0}G(x,t;L,0)
&= \sum_{\substack{n=-\infty \\ (n\ne 0)}}^{\infty}
   \frac{-i\sin n\pi x/L\,e^{i(n\pi ct/L)}}{cL\cos n\pi} \notag\\
&= \frac{2}{cL}\sum_{n=1}^{\infty}(-1)^n\sin\frac{n\pi x}{L}\sin\frac{n\pi ct}{L},
\label{eq:13.8.9}
\end{align}
where the positive and negative $n$ contributions have been combined into one term.
Finally, using \eqref{eq:13.8.4}, we obtain
\[
u(x,t) = \sum_{n=1}^{\infty} A_n(t)\sin\frac{n\pi x}{L},
\quad\text{where}\quad
A_n(t) = -(-1)^n\frac{2c}{L}\int_0^t b(t_0)\sin\frac{n\pi c}{L}(t-t_0)\,dt_0,
\]
the same result as would be obtained by the method of eigenfunction expansion.

The influence function is an infinite series of the eigenfunctions.  For homogeneous problems with homogeneous boundary conditions, inverting the Laplace transform using an infinite sequence of poles also will yield a series of eigenfunctions, the same result as obtained by separation of variables.  In fact, it is the Laplace transform method that is often used to \emph{prove} the validity of the method of separation of variables.

%% ── Exercises 13.8 ──────────────────────────────────────────────
\begin{exercises}{13.8}

\item[\starred 13.8.1.] Solve for $u(x,t)$ using Laplace transforms
\[
\pdd{u}{t} = c^2\pdd{u}{x}
\]
\[
u(x,0) = f(x) \qquad \pd{u}{t}(x,0) = 0 \qquad u(0,t) = 0 \qquad u(L,t) = 0.
\]
Invert the Laplace transform of $u(x,t)$ using the residue theorem for contour integrals in the $s$-plane.  Show that this yields the same result as derivable by separation of variables.

\item Modify Exercise~13.8.1 if instead
\begin{enumerate}
\item[(a)] $u(x,0) = 0$ \quad and\quad $\pd{u}{t}(x,0) = g(x)$
\item[(b)] $u(0,t) = 0$ \quad and\quad $\pd{u}{x}(L,t) = 0$
\item[(c)] $\pd{u}{x}(0,t) = 0$ \quad and\quad $\pd{u}{x}(L,t) = 0$
\end{enumerate}

\item Solve for $u(x,t)$ using Laplace transforms
\[
\pd{u}{t} = k\pdd{u}{x}
\]
subject to $u(x,0) = f(x)$, $u(0,t) = 0$, and $u(L,t) = 0$.
Invert the Laplace transform of $u(x,t)$ using the residue theorem for contour integrals in the complex $s$-plane.  By what other method can this representation of the solution be obtained?  (Compare to Exercise~13.5.6.)

\item Consider
\[
\pdd{u}{t} = c^2\pdd{u}{x} + \sin\sigma_0 t
\]
\[
u(x,0) = 0 \qquad u(0,t) = 0
\]
\[
\pd{u}{t}(x,0) = 0 \qquad u(L,t) = 0.
\]
\begin{enumerate}
\item[(a)] Solve using Laplace transforms (with contour inversion) if $\sigma_0 \ne c(m\pi/L)$.
\item[(b)] Solve if $\sigma_0 = c(3\pi/L)$.  Show that resonance occurs (see Sec.~8.5).
\end{enumerate}

\end{exercises}
